\section{浮点数表示基础}

\subsection{FP32：标准单精度浮点数}

在深度学习中，模型参数和梯度通常以浮点数的形式存储在GPU显存中。FP32（32位浮点数）是最基本的数据类型，每个参数占用 \textbf{4字节}。

\begin{figure}[H]
\centering
\includegraphics[width=0.95\textwidth]{slides-images/slide-003.jpg}
\caption{FP32 浮点数格式：1位符号位 + 8位指数位 + 23位尾数位，共32位（4字节）\protect\footnotemark}
\end{figure}
\footnotetext{Slides 第4页。}

FP32 的数值由以下公式计算：

$$(-1)^{B} \times 2^{E-127} \times \left(1 + \sum_{i=1}^{23} b_{23-i} \cdot 2^{-i}\right)$$

\begin{itemize}
    \item $B$：符号位（sign），决定正负
    \item $E$：指数位（exponent，8位），决定\textbf{动态范围}——能表示多大或多小的数
    \item $b_i$：尾数位（mantissa，23位），决定\textbf{精度}——数值的精确程度
\end{itemize}

\begin{figure}[H]
\centering
\includegraphics[width=0.95\textwidth]{slides-images/slide-005.jpg}
\caption{FP32 数值计算公式：符号位、指数位、尾数位分别对应正负号、范围、精度\protect\footnotemark}
\end{figure}
\footnotetext{Slides 第6页。}

\begin{knowledgebox}{指数位与尾数位的权衡}
\textbf{指数位越多}，能表示的数的范围越大（更小的数和更大的数都能表示）。\\
\textbf{尾数位越多}，数值的精度越高（相邻可表示数之间的间隔更小）。\\
在总位数固定的前提下，指数位与尾数位之间存在权衡——这是理解FP16、BF16等格式的关键。
\end{knowledgebox}

\subsection{FP16：半精度浮点数}

FP16 只有16位（2字节），内存需求是 FP32 的一半。其格式为：1位符号位 + 5位指数位 + 10位尾数位。

\begin{figure}[H]
\centering
\includegraphics[width=0.95\textwidth]{slides-images/slide-008.jpg}
\caption{FP16 与 FP32 对比：FP16 指数位仅5位（动态范围小），尾数位10位（精度低）\protect\footnotemark}
\end{figure}
\footnotetext{Slides 第9页。}

\begin{table}[H]
\centering
\begin{tabular}{lccc}
\toprule
\textbf{数据类型} & \textbf{总位数} & \textbf{指数位} & \textbf{尾数位} \\
\midrule
FP32  & 32 & 8  & 23 \\
FP16  & 16 & 5  & 10 \\
BF16  & 16 & 8  & 7  \\
\bottomrule
\end{tabular}
\caption{三种常见浮点数格式的位数分配}
\end{table}

FP16 的两个核心问题：

\begin{enumerate}
    \item \textbf{动态范围不足}：指数位仅5位，导致非常小的数会被\textbf{下溢为零}（underflow）。例如，小于约 $6 \times 10^{-5}$ 的数在 FP16 中直接变为0。
    \item \textbf{精度不足}：尾数位仅10位，导致\textbf{舍入误差}。例如，1.0001 在 FP16 中会被舍入为 1.0。
\end{enumerate}

\begin{figure}[H]
\centering
\includegraphics[width=0.95\textwidth]{slides-images/slide-010.jpg}
\caption{FP16 下的梯度下溢问题：训练过程中的激活梯度分布，超过一半的梯度在 FP16 中会变为零\protect\footnotemark}
\end{figure}
\footnotetext{Slides 第11页。图来自 NVIDIA 博客。}

\begin{warningbox}{FP16 梯度下溢是致命问题}
在神经网络训练中，梯度通常非常小。如果直接用 FP16 计算梯度，\textbf{超过一半的梯度值会被四舍五入为零}，导致模型无法有效学习。这就是为什么不能简单地将所有计算切换到 FP16。
\end{warningbox}

\subsection{本章小结}

浮点数的位数分配（指数位 vs 尾数位）决定了动态范围和精度之间的权衡。FP32 是标准格式（4字节/参数），FP16 节省一半内存但存在严重的动态范围和精度问题。理解这些限制是混合精度训练的前提。

%% ========== 混合精度训练 ==========

